Decision trees~\cite{breiman1984cart}, rule lists~\cite{letham2015stroke, angelino2018learningcertifiablyoptimalrule} and prototype-based classifiers~\cite{chen2019looks, bien2012prototype} are typical interpretable-by-design models; among them, GAMs~\cite{hastie1990generalized, wood2017generalized} are popular for their simple global interpretability.

GAMs model the output as $g(\mathbb{E}[y|\mathbf{x}]) = \beta_0 + \sum_{i} f_i(x_i)$, with univariate shape functions $f_i$, intercept $\beta_0$ and link function $g$.
They differ mainly in how $f_i$ are learned: penalized splines~\cite{wood2017generalized} give smooth curves, boosted bagged shallow trees~\cite{lou2012intelligible, nori2019interpretml} give step functions that follow sharp changes, and neural variants learn one network per feature~\cite{agarwal2021neural}, a set of basis functions shared by all features~\cite{radenovic2022neural}, or a boosted ensemble of extreme learning machines~\cite{KRAUS2024IGANN}.

\gamms~add pairwise terms, $g(\mathbb{E}[y|\mathbf{x}]) = \beta_0 + \sum_{i} f_i(x_i) + \sum_{(i,j)\in\mathcal{I}} f_{ij}(x_i, x_j)$, over a set $\mathcal{I}$ of selected pairs, and differ in how they choose $\mathcal{I}$ and fit the terms.
GA$^2$M~\cite{lou2013accurate} fits boosted trees and adds pairs greedily, ranked by how much each reduces the residual loss;
GAMI-Net~\cite{yang2020gami} fits one network per term and considers only pairs with at least one already-selected main effect;
GRAND-SLAMIN~\cite{ibrahim2023grandslamin} fits soft trees and selects pairs end to end through differentiable gates, optionally under hierarchy constraints;
and NODE-GA$^2$M~\cite{chang2022node} fits differentiable oblivious trees, each restricted to at most two features.

CALM draws inspiration from \emph{regional effect methods}, post-hoc explainability methods that extend \emph{global effect plots}.
Global effect plots, such as Partial Dependence Plots (PDP)~\cite{friedman2001greedy}, Accumulated Local Effects (ALE)~\cite{apley2020visualizing} and SHAP Dependence Plots (SHAP-DP)~\cite{lundberg2020local}, explain a black-box model $f$ by decomposing it into $d$ one-dimensional plots, one per feature $x_i$.
When $f$ contains strong interactions, these plots average heterogeneous local effects into one curve and can mislead, a phenomenon known as aggregation bias~\cite{mehrabi2021survey, herbinger2022REPID, baniecki2024robustness, gkolemis2023rhale}.

Regional effect methods address this by splitting the input space with a tree into subregions where interactions are weak, and drawing each feature's plot separately in every subregion.
REPID~\cite{herbinger2022REPID} grows one tree per feature, while GADGET~\cite{herbinger2024GADGET}, FD-Tree~\cite{laberge2024fdtree} and GRANITE~\cite{herbinger2026granite} grow one tree shared by several features.
CALM adopts GADGET's heterogeneity criterion and REPID's per-feature trees to find its regions (Step~1(b), Section~\ref{subsec:training}).
Per-feature trees let each feature split only on the features it interacts with, whereas a shared tree imposes the same splits on every feature.
In Figure~\ref{fig:concept-image}, a shared tree would split the humidity, which interacts with no feature, on the working day only because the hour needs it.
Finally, we clarify that GADGET, REPID and all regional effect methods are not competitors of CALM: they explain a fitted model post hoc, whereas a CALM is itself a predictor.

CALM also relates to model distillation~\cite{hinton2015distilling, tan2018distill} and surrogate models~\cite{guidotti2018survey}, which train an interpretable student to mimic a black-box teacher, locally~\cite{ribeiro2016should} or globally~\cite{guidotti2018survey}.
However, instead of explaining the teacher, a CALM uses it only to find where interactions are weak, and then fits its curves to the data.
In a rough analogy, CALM serves as the interpretable-by-design student that replaces the black-box teacher as the final predictor.
SLIM~\cite{hu2020surrogate} and other model-based trees~\cite{herbinger2023leveraging} fit a simple model in each leaf of a single tree, mainly as surrogates that explain a black box, so all features share one partition.
Regionally additive models~\cite{gkolemis2023regionally} also build an interpretable model from regional effects, but do not characterize its interpretability or evaluate it systematically.
